% BEGIN JE20261104_EXTENSION
\paragraph{Corollary 1 (the fee-adjustment sign contrast).}\label{cor:repricing}
Under the strict primitive-share band in Proposition~1, fix any reference
$\lambda_0$ in its compact feasible interval. Consider a diagnostic
counterfactual that holds the private-only posted fee at the optimal
$f_i^P(\lambda_0)$ while varying $\lambda$ locally; the free-basic-plus-premium
continuation is otherwise unchanged. Let $T_i^{\mathrm{fr}}$ be its gross
conditional public cost ceiling, and $T_i$ the equilibrium-fee ceiling.
For all sufficiently weak positive matching intensities,
\begin{equation}
 \left.\frac{dT_i^{\mathrm{fr}}}{d\lambda}\right|_{\lambda_0}<0
 <\left.\frac{dT_i}{d\lambda}\right|_{\lambda_0}.
 \label{eq:repricing-reversal}
\end{equation}
Indeed, with $D=d+\eta_i\lambda$, freezing the fee locally gives
$\ell^{\mathrm{fr}}(\lambda)=
[f_i^P(\lambda_0)-U]/D$ and hence
$(\ell^{\mathrm{fr}})'|_{\lambda_0}=-\eta_i\ell(\lambda_0)/D(\lambda_0)<0$.
The first-order fee-freezing effect and the already negative
$K_i$ together make the cost-ceiling slope negative;
reoptimisation instead increases exclusion, yielding Proposition~1's
positive slope. This is a mechanism-isolating \emph{off-equilibrium fee}
comparison, not a second pricing equilibrium or a policy instrument.
The proof is in Online Resource~1.

\paragraph{Proposition 2 (regular-distribution screening--incidence boundary).}
\label{prop:general-distribution}
The conditional sign inversion is not confined to uniformly distributed
skills. Replace the local skill distribution by a continuously distributed
$F_i$ on $[0,1]$, with positive twice continuously differentiable density
$g_i$ on the interior and nondecreasing hazard
$h_i(t)=g_i(t)/[1-F_i(t)]$. Retain $z_i=\theta_i$ and
all other technological, payoff and timing primitives; hold the other
region at full participation with mean skill
$\bar\theta_j=\mathbb E[\theta_j]>0$.
Assume $u_i h_i(0+)<d=\Delta_i$, positive baseline $p_b$,
$M_i<0$, $\Pi>0$, and a compact feasible $\lambda$ interval.
Let $\ell_0\in(0,1)$ be the unique global private-only admission cutoff
at zero matching intensity, satisfying
$d[1-F_i(\ell_0)]=(u_i+d\ell_0)g_i(\ell_0)$. Write
\begin{equation}
 \begin{gathered}
 J_F=\int_0^{\ell_0}t g_i(t)\,dt,\quad
 B_F=2d g_i(\ell_0)+(u_i+d\ell_0)g_i'(\ell_0)>0,\\
 A_F=\frac{u_i(u_i+d\ell_0)g_i(\ell_0)^2}{dB_F}>0,
 \qquad
 \omega_{i,F}^{\mathrm{crit}}
   =\frac{-M_i}{\Pi}\frac{J_F}{J_F+A_F}.
 \end{gathered}\label{eq:general-share-band}
\end{equation}
For sufficiently small positive $\gamma=\rho\mu$, uniformly on that
$\lambda$ interval, the gross conditional public-access ceiling obeys
\begin{equation}
 \frac{dT_i^F}{d\lambda}
 =\gamma N_iN_j\bar\theta_j
   \bigl[M_i J_F+\omega_i\Pi(J_F+A_F)\bigr]
   +O(\gamma^2).
 \label{eq:general-slope}
\end{equation}
Consequently, if the feasible share band
$\omega_{i,F}^{\mathrm{crit}}<\omega_i<\min\{1,-M_i/\Pi\}$
is nonempty, then $K_i<0$, $S_i^0<0$, and
$dT_i^F/d\lambda>0$ throughout the compact interval for sufficiently
small positive $\gamma$. An \emph{additional} positive level of the
cost ceiling requires that the zero-matching direct-service gain
$N_i\{u_iF_i(\ell_0)+d[q_i(\tau_B)-q_i(\ell_0)]\}$ be strictly
positive, where $q_i(t)=\int_t^1xg_i(x)\,dx$ and $\tau_B$ is the
optimal co-provision premium cutoff. For the uniform distribution and
$\bar\theta_j=1/2$, equation~\eqref{eq:general-slope} reduces exactly
to Proposition~1. This is a conditional \emph{distributional extension},
not a theorem for arbitrary skill--taste correlation, endogenous partner
coverage or the full four-entrant game. Proof and the sufficient
hazard-based global fee condition are given in Online Resource~1.

\paragraph{Proposition 3 (private-screening statistics do not determine the policy-value slope).}
\label{prop:screening-insufficient}
Even within the original uniform-ability Cournot model, identical private
fees, founder admission, selected abilities, success probabilities and
venture operating profits need not identify the sign of the conditional
regional public-access cost-ceiling derivative.
To see this constructively, retain every original primitive except the
two net demand intercepts in region $B$. Keep $n_B=2$,
$b_B=d_B=1$, $\chi_B=9/10$, and the successful venture's output
$y_B=20/73$. For $x\in\{55/219,\,2/5\}$ set the incumbent firm's
Cournot output to $x$, with net demand intercepts
\begin{equation}
 A_B(x)=3x+\tfrac9{10}\tfrac{20}{73},\qquad
 D_B(x)=2\tfrac{20}{73}+\tfrac95 x.
 \label{eq:incidence-twin-cournot}
\end{equation}
Both specifications have strictly interior Cournot outputs and the
\emph{same} total venture profit $\Pi$, hence the same
founder matching returns and private-only fee response at each
$\lambda\in[0,1/2]$ when region $A$ fully participates.
Nevertheless, both have $K_B<0$ and $S_B^0<0$, while
\begin{equation}
 \left.\frac{dT_B}{d\lambda}\right|_{x=55/219}
    >\tfrac1{10000}>0,\qquad
 \left.\frac{dT_B}{d\lambda}\right|_{x=2/5}
    <-\tfrac1{100000}<0
 \quad\text{for every }\lambda\in[0,1/2].
 \label{eq:incidence-twin-signs}
\end{equation}
The distinction is solely in consumer and incumbent incidence
$M_B(x)=-4(219x-14)/5329$. Thus the private service fee,
participation statistics and venture profit are \emph{not sufficient
statistics} for the local policy-value comparison: downstream market
incidence must also be evaluated. The construction does not claim that
the governments' complete entry-game equilibria coincide.
A fully rational, global-in-$\lambda$ proof is in Online Resource~1.

% END JE20261104_EXTENSION
